Land-surface schemes regulate the exchange of momentum, energy, water, and
carbon between the atmosphere and the terrestrial surface. Their treatment of
soil moisture, snow, vegetation, and surface temperature influences
near-surface meteorology, boundary-layer development, runoff, and the memory
of land--atmosphere anomalies.

Land-surface modelling rests on several decades of process development
accumulated in a small number of community schemes. The Common Land Model,
which assembled the land physics of earlier biosphere--atmosphere schemes
into a single column model for energy, water, and carbon exchange
\citep{Dai2003CoLM}, is widely used in climate and Earth system modelling
and has been carried forward along two main lines: the Community Land Model
(CLM) developed in the United States \citep{Oleson2010CLM} and the Common
Land Model (CoLM) developed in China, most recently documented as CoLM 2024
\citep{Dai2004TwoBigLeaf,YuanDai2025CoLM}. The Noah scheme and its
multi-parameterization successor Noah-MP carry the land surface of mesoscale
numerical weather prediction \citep{Niu2011NoahMP}. Operational
numerical-weather-prediction centres have likewise long developed their own
land-surface models, among them ECMWF's ECLand, the Joint UK Land
Environment Simulator (JULES; \citealp{Wiltshire2020JULES}), and the
Community Atmosphere--Biosphere Land Exchange model (CABLE;
\citealp{Haverd2018CABLE}). At ECMWF, in particular, the land surface evolved from
the original TESSEL formulation through the tiled HTESSEL hydrology
\citep{ViterboBeljaars1995,VanDenHurk2000,Balsamo2009} into ECLand, a
modular modelling system that couples energy, water, and carbon and
underpins the land surface of the operational IFS and its ensemble
reanalyses \citep{Boussetta2021ECLand}. The OpenIFS land-surface component
examined here belongs to this lineage; its physical formulation is
documented in the IFS physical-process documentation \citep{ECMWFIFSPartIV}.
These schemes are mature in their process content, but they are written in
Fortran, and neither derivatives nor accelerator-oriented batch execution is
part of their design, so parameter and sensitivity studies rely on ensembles
of independent integrations.

Atmospheric prediction has in the same years moved toward models whose
complete time integration is differentiable. Differentiable programming was
brought into Earth-system modelling to combine process models with learned
components inside one gradient computation \citep{Gelbrecht2023}, and reviews
of differentiable geoscience identify gradient access through the process
equations as the requirement that unifies physical and machine-learning
approaches \citep{Shen2023}. NeuralGCM couples a differentiable dynamical
core with learned parameterizations and trains the combination end to end
\citep{Kochkov2024NeuralGCM}; its precipitation-focused successor is trained
directly against satellite retrievals, which is only possible because every
component of the model can be differentiated \citep{Yuval2026Precip}. These
developments show what differentiability offers operational prediction: the
observation record itself becomes the training signal.

Differentiable land-surface implementations are more recent and remain
site-focused. JAX-CanVeg reimplements the CanVeg canopy model in JAX, keeps
the ecohydrological process chain intact, and embeds a neural-network
stomatal-conductance closure within the process equations
\citep{Jiang2025JAXCanVeg}. NoahPy recasts the Noah land surface model, in a
permafrost thermo-hydrology configuration, as a process-encapsulated
recurrent structure in PyTorch; it reproduces the Fortran original at site
level with Nash--Sutcliffe efficiency above 0.99, and gradient-based
calibration of its soil parameters converges faster and more stably than the
SCE-UA search at a Qinghai--Tibet Plateau site \citep{NoahPy2026}. No
differentiable implementation yet covers the complete land physics of an
operational forecasting system over a continental domain with decadal
verification, and the conditioning of the resulting estimation problem,
which decides whether a gradient-based fit is well posed, has received
little attention in land-surface applications.

Land-surface simulations now appear in applications that require large
ensembles, parameter sensitivities, or the assimilation of land-surface
observations. These applications benefit from an implementation that
evaluates many independent columns concurrently and provides derivatives of
selected model outputs with respect to inputs and parameters.
Finite-difference ensembles need one integration per perturbed control;
automatic differentiation obtains every selected derivative from a single
reverse calculation while retaining process-based state variables and
conservation constraints.

This study presents a PyTorch implementation of the OpenIFS SURF land-surface
physics. The implementation retains the tiled process sequence of the reference scheme:
surface exchange, the implicit surface energy balance, canopy interception,
soil heat and water transport, and single-layer snow processes. Every
conditional regime switch of the reference scheme is kept unsmoothed, so the
operator that is differentiated is the same operator whose numerical
equivalence is verified. Its purpose is
to reproduce the numerical behaviour of SURF while enabling batched CPU/GPU
execution and differentiation of the tensor operations. The implementation is
therefore evaluated as a numerical translation of an established physical
scheme rather than as a new land-surface parameterization. Verification and
estimation are given equal weight: equivalence against the Fortran reference
is established before any gradient is used.

The study has six objectives. First, the principal SURF equations and their
tensor representation are described. Second, the implementation is evaluated
against the independent Fortran reference implementation under identical
roughness and surface-layer exchange boundary conditions, a separate native OpenIFS offline-column experiment
isolates differences associated with reconstructing the atmospheric
surface-layer boundary in standalone mode, and a one-year full-domain
autonomous integration over heterogeneous CMFD forcing tests the complete
prognostic sequence from a common initial state. Third, a ten-year successively restarted
integration examines the stability of this equivalence over multiple forcing
years. Fourth,
multi-step automatic derivatives of forcing, initial state, and vegetation
controls are verified against centred finite differences and applied in
regional sensitivity analysis, observation-space identifiability diagnosis,
parameter calibration, and initial-state retrieval. Fifth, the autonomous
integration is evaluated against independent Terra MODIS observations of land
surface temperature and snow cover. Sixth, and central to the
scientific aim, the differentiability is used to close the loop between the
MODIS observations and the process formulation: the residual daytime bias is
separated into a parameter-explained part and a structural remainder, the
missing process is diagnosed by reverse-mode attribution, and a new
finite-capacity surface-layer closure is designed, calibrated, and verified
against the same observations. The batch-size CPU
and GPU experiments that quantify the acceleration of the batched tensor
implementation close the study, without changing the physical time step or
numerical precision.

Section~\ref{sec:model-background} describes the model formulation.
Section~\ref{sec:fortran-torch-equivalence} presents the Fortran--PyTorch
consistency experiments, including the ten-year successively restarted integration of
Sect.~\ref{sec:multiyear}. Section~\ref{sec:gradient-calibration} describes
the differentiable formulation and its use in sensitivity analysis and
initial-state retrieval. Section~\ref{sec:modis-evaluation} presents the
independent MODIS evaluation and the structural daytime bias it reveals, then
uses the differentiable formulation to attribute that bias to a missing
near-surface process and to design, calibrate, and verify a finite-capacity
surface-layer closure that removes most of it; the
parameter-calibration and pseudo-forcing experiments that establish the
structural nature of the bias and motivate the closure are documented in
Appendix~\ref{app:modis-calibration}. Section~\ref{sec:performance}
reports computational scaling.
